\documentclass[11pt]{article}

\usepackage[T1]{fontenc}
\usepackage{lmodern}
\usepackage[margin=1in]{geometry}
\usepackage{amsmath,amssymb,amsthm,mathtools}
\usepackage{microtype}
\usepackage[hidelinks]{hyperref}

\newtheorem{theorem}{Theorem}
\newtheorem{lemma}[theorem]{Lemma}
\newtheorem{corollary}[theorem]{Corollary}
\newtheorem{remark}[theorem]{Remark}

\newcommand{\ind}{\mathbf{1}}
\newcommand{\cost}{\operatorname{cost}}
\newcommand{\OPT}{\operatorname{OPT}}
\newcommand{\ALG}{\operatorname{ALG}}

\allowdisplaybreaks

\hypersetup{
  pdftitle={A Deterministic 18-Competitive Algorithm for Online Consensus Clustering},
  pdfsubject={Online consensus clustering with logarithmic per-vertex recourse}
}

\begin{document}

\section{A Deterministic 18-Competitive Algorithm for Online Consensus Clustering}

\subsection{Problem definition and notation}

Let $V=\{v_1,\ldots,v_n\}$ be the vertex set, where the vertices arrive
online in the order $v_1,v_2,\ldots,v_n$.  At time $t$, the set of vertices
that have arrived is
\[
  V_t=\{v_1,\ldots,v_t\}.
\]
We are given $m$ input clusterings (equivalently, partitions)
\[
  \mathcal P^{(1)},\mathcal P^{(2)},\ldots,\mathcal P^{(m)}
\]
of $V$.  When a vertex arrives, its cluster in each of the $m$ input
clusterings is revealed.

For a statement $E$, let
\[
  \ind[E]=
  \begin{cases}
    1,&\text{if $E$ is true},\\
    0,&\text{if $E$ is false}
  \end{cases}
\]
denote its indicator.  For two distinct vertices $u,v\in V$, define
\[
 d_{uv}
 =\frac{1}{m}\sum_{i=1}^m
   \ind[\mathcal P^{(i)}\text{ separates }u\text{ and }v].
\]
Thus, $d_{uv}$ is the fraction of input clusterings that place $u$ and $v$
in different clusters, and $1-d_{uv}$ is the fraction that place them in the
same cluster.  Notice that $0\le d_{uv}\le1$.  Moreover, $d$ is a
semimetric: it is symmetric, $d_{uu}=0$, and
\[
  d_{uv}\le d_{uw}+d_{wv}
\]
for all vertices $u,v,w$.  Indeed, the separation indicator of every input
partition satisfies the triangle inequality, and averaging preserves it.

For an output clustering $\mathcal C$ of $V_t$, write
$u\sim_{\mathcal C}v$ when $u$ and $v$ belong to the same output cluster,
and write $u\not\sim_{\mathcal C}v$ otherwise.  The normalized consensus
cost of $\mathcal C$ at time $t$ is
\[
 \cost_t(\mathcal C)
 =\sum_{\{u,v\}\subseteq V_t}
   \left(
     d_{uv}\ind[u\sim_{\mathcal C}v]
     +(1-d_{uv})\ind[u\not\sim_{\mathcal C}v]
   \right).
\]
Every sum indexed by $\{u,v\}\subseteq V_t$ is over unordered pairs of
distinct vertices.  Define the pair-disagreement distance between two
clusterings $\mathcal C$ and $\mathcal P$ of $V_t$ by
\[
 \operatorname{dist}_t(\mathcal C,\mathcal P)
 =\sum_{\{u,v\}\subseteq V_t}
   \ind\!\left[
     \ind[u\sim_{\mathcal C}v]
     \ne
     \ind[u\sim_{\mathcal P}v]
   \right].
\]
The usual consensus objective
$\sum_{i=1}^m \operatorname{dist}_t(\mathcal C,\mathcal P^{(i)})$
equals $m\cost_t(\mathcal C)$, so normalization does not change any
approximation or competitive ratio.  Let
\[
  \OPT_t=\min_{\mathcal C}\cost_t(\mathcal C),
\]
where the minimum is over all clusterings of $V_t$.

For a real number $x$, its \emph{positive part} is
\[
  (x)_+ := \max\{x,0\}.
\]
In particular,
\[
 (1-2d_{uv})_+
 =
 \begin{cases}
  1-2d_{uv},&\text{if }d_{uv}\le \tfrac12,\\
  0,&\text{if }d_{uv}>\tfrac12.
 \end{cases}
\]
The subscript ``$+$'' therefore does not denote addition; it means that a
negative value is replaced by zero.

For every unordered pair $\{u,v\}$, define
\begin{align*}
 \ell_{uv}&=\min\{d_{uv},1-d_{uv}\},\\
 p_{uv}&=(1-2d_{uv})_+,\\
 n_{uv}&=(2d_{uv}-1)_+.
\end{align*}
Here $p_{uv}$ is the confidence in joining $u$ and $v$, $n_{uv}$ is the
confidence in separating them, and $\ell_{uv}$ is the unavoidable
pairwise ambiguity.  Directly from the definitions,
\begin{equation}
 d_{uv}=\ell_{uv}+n_{uv},
 \qquad
 1-d_{uv}=\ell_{uv}+p_{uv},
 \qquad
 p_{uv}+n_{uv}+2\ell_{uv}=1.
 \label{eq:pair-identities}
\end{equation}

For a finite set $A$, $|A|$ denotes its number of elements.  For two
disjoint nonempty vertex sets $A,B$, define
\begin{align*}
 D(A,B)&=\sum_{u\in A}\sum_{v\in B}d_{uv},\\
 L(A,B)&=\sum_{u\in A}\sum_{v\in B}\ell_{uv},\\
 P(A,B)&=\sum_{u\in A}\sum_{v\in B}p_{uv},\\
 N(A,B)&=\sum_{u\in A}\sum_{v\in B}n_{uv}.
\end{align*}
The average distance between $A$ and $B$ is
\[
  \overline d(A,B)=\frac{D(A,B)}{|A||B|}.
\]
For a vertex $z\notin C$, the notation $D(z,C)$ means
$D(\{z\},C)$, and similarly for $L(z,C),P(z,C)$, and $N(z,C)$.
Summing~\eqref{eq:pair-identities} over $A\times B$ gives
\begin{equation}
 D(A,B)=L(A,B)+N(A,B),
 \qquad
 P(A,B)-N(A,B)=|A||B|-2D(A,B).
 \label{eq:set-identities}
\end{equation}

Finally, define
\[
 L_t=\sum_{\{u,v\}\subseteq V_t}\ell_{uv}.
\]
Every clustering pays either $d_{uv}$ or $1-d_{uv}$ on pair $\{u,v\}$,
and both quantities are at least $\ell_{uv}$.  Consequently,
\begin{equation}
  L_t\le \OPT_t.
  \label{eq:ambiguity-lower-bound}
\end{equation}

\paragraph{Recourse convention.}
Every output cluster carries a persistent identifier.  An old vertex incurs
one unit of recourse whenever its cluster identifier changes between two
consecutive time steps.  The initial identifier assigned to a newly arriving
vertex is not counted as recourse.  This is the standard cluster-identifier
definition of per-vertex recourse.

\subsection{The algorithm}

At every time $t$, the algorithm maintains a clustering of $V_t$.

\begin{enumerate}
 \item When $v_t$ arrives, create the singleton cluster $\{v_t\}$ with a
       fresh identifier.
 \item While there are two current clusters $A,B$ satisfying
       \begin{equation}
         \overline d(A,B)=\frac{D(A,B)}{|A||B|}<\frac14,
         \label{eq:merge-rule}
       \end{equation}
       merge $A$ and $B$.
 \item If several pairs satisfy~\eqref{eq:merge-rule}, choose one using any
       fixed deterministic tie-breaking rule.
 \item The merged cluster inherits the identifier of its larger child.  If
       $|A|=|B|$, use a fixed deterministic tie-breaking rule to choose the
       inherited identifier.
\end{enumerate}

A cluster is called \emph{reachable} if it can be constructed from singleton
clusters through a sequence of merges satisfying~\eqref{eq:merge-rule}.
Every cluster maintained by the algorithm is reachable.  The sequence of
merges constructing a reachable cluster can be represented by a binary
\emph{merge tree}: its leaves are singleton vertices, and every internal node
is the union of its two children.

\begin{theorem}[Main theorem]
For every time $t$ and every arrival order, the above deterministic algorithm
satisfies
\[
  \cost_t(\mathcal C_t)\le18\OPT_t,
\]
where $\mathcal C_t$ is the clustering output at time $t$.  Moreover, every
vertex changes its cluster identifier at most $\lfloor\log_2 n\rfloor$ times.
The competitive ratio and recourse bound are independent of $m$.
\end{theorem}

Here $\log_2 n$ is the base-$2$ logarithm, and $\lfloor x\rfloor$ denotes
the greatest integer at most $x$.

\subsection{Analysis of the algorithm}

\paragraph{High-level idea of the analysis.}
We first give an informal overview of the proof.  Recall that
\[
    \ell_{uv}=\min\{d_{uv},1-d_{uv}\}.
\]
Regardless of whether a clustering joins or separates $u$ and $v$, it pays
at least $\ell_{uv}$ on this pair.  Therefore,
\[
    L_t=\sum_{\{u,v\}\subseteq V_t}\ell_{uv}\le \OPT_t.
\]
Our goal is to prove that the cost of the algorithm is at most a constant
times $L_t$.

The quantities
\[
    p_{uv}=(1-2d_{uv})_+
    \qquad\text{and}\qquad
    n_{uv}=(2d_{uv}-1)_+
\]
measure the majority margins in favor of joining and separating $u$ and
$v$, respectively.  In particular,
\[
    d_{uv}=\ell_{uv}+n_{uv},
    \qquad
    1-d_{uv}=\ell_{uv}+p_{uv}.
\]
Thus, if the algorithm joins $u$ and $v$, then the additional cost above
the unavoidable ambiguity $\ell_{uv}$ is $n_{uv}$.  If the algorithm
separates them, then the additional cost is $p_{uv}$.

The main structural goal is to prove that, for any two clusters $A$ and
$B$ that can be produced by the algorithm,
\begin{equation}
    \min\{P(A,B),N(A,B)\}\le 5L(A,B).
    \label{eq:informal-pair-coherence}
\end{equation}
In words, between two reachable clusters, the smaller of the total joining
confidence and the total separating confidence is bounded by five times the
total ambiguity.

We establish~\eqref{eq:informal-pair-coherence} in three steps.

\begin{enumerate}
    \item
    Lemma~2, called \emph{radial coherence}, considers one vertex $z$ and
    one reachable cluster $C$.  It proves
    \[
        \min\{P(z,C),N(z,C)\}\le 2L(z,C).
    \]
    Thus, when a reachable cluster is viewed from any external vertex, it
    cannot simultaneously contain a large amount of evidence for joining
    and a large amount of evidence for separating, unless there is also a
    sufficiently large amount of ambiguity.

    The proof is by induction on the merge tree of $C$.  The only difficult
    case is when one child of a merge is predominantly close to $z$, while
    the other child is predominantly far from $z$.  We show that such two
    children cannot have average distance less than $1/4$, and therefore the
    algorithm could not have merged them.

    \item
    Lemma~3, called the \emph{polarization matrix lemma}, is an algebraic
    local-to-global statement.  Consider a matrix $(q_{ij})$ with
    \[
        p_{ij}=(q_{ij})_+,\qquad
        n_{ij}=(-q_{ij})_+,\qquad
        u_{ij}=1-|q_{ij}|.
    \]
    If every row and every column satisfies
    \[
        \min\{P_\lambda,N_\lambda\}\le U_\lambda,
    \]
    then the whole matrix satisfies
    \[
        \min\left\{\sum_{i,j}p_{ij},\sum_{i,j}n_{ij}\right\}
        \le \frac52\sum_{i,j}u_{ij}.
    \]
    Thus, row-wise and column-wise coherence imply coherence for the entire
    matrix.

    \item
    Corollary~4 applies Lemma~3 to the matrix
    \[
        q_{uv}=1-2d_{uv},
        \qquad u\in A,\ v\in B.
    \]
    Lemma~2 verifies the required inequality for every row because $B$ is
    reachable, and for every column because $A$ is reachable.  Since
    \[
        1-|q_{uv}|=1-|1-2d_{uv}|=2\ell_{uv},
    \]
    Lemma~3 yields
    \[
        \min\{P(A,B),N(A,B)\}
        \le \frac52\cdot 2L(A,B)
        =5L(A,B),
    \]
    which is exactly~\eqref{eq:informal-pair-coherence}.
\end{enumerate}

We then use Corollary~4 to analyze the cost of the algorithm.  There are two
types of pairs.

First, suppose that $A$ and $B$ are the two children of an internal node of
a merge tree.  Since the algorithm merged them,
\[
    D(A,B)<\frac{|A||B|}{4}.
\]
Consequently,
\[
    P(A,B)-N(A,B)
    =|A||B|-2D(A,B)
    >0.
\]
Hence $N(A,B)$ is the smaller quantity in
$\{P(A,B),N(A,B)\}$.  Corollary~4 gives
\[
    N(A,B)\le 5L(A,B).
\]
The algorithm joins all pairs in $A\times B$, and therefore pays
\[
    D(A,B)=L(A,B)+N(A,B)\le 6L(A,B).
\]

Second, suppose that $A$ and $B$ are two distinct terminal clusters.  Since
the merge loop has terminated,
\[
    D(A,B)\ge\frac{|A||B|}{4}.
\]
The algorithm separates all pairs in $A\times B$.  If
$P(A,B)\le N(A,B)$, Corollary~4 immediately gives
\[
    L(A,B)+P(A,B)\le 6L(A,B).
\]
Otherwise, Corollary~4 gives $N(A,B)\le 5L(A,B)$, and hence
\[
    D(A,B)=L(A,B)+N(A,B)\le 6L(A,B).
\]
Terminality gives $|A||B|\le 4D(A,B)$, so the separated cost is
\[
    |A||B|-D(A,B)
    \le 3D(A,B)
    \le 18L(A,B).
\]

The internal merge-tree cuts and the pairs of distinct terminal clusters
partition all unordered pairs of vertices.  Therefore,
\[
    \cost_t(\mathcal C_t)
    \le 18L_t
    \le 18\OPT_t.
\]
Finally, whenever an old vertex changes its cluster identifier, its current
cluster is the smaller child of a merge, and hence its cluster size at least
doubles.  Therefore, every vertex changes its identifier at most
$\lfloor\log_2 n\rfloor$ times.

We now prove the three structural statements used above.

\begin{lemma}[Radial coherence]
\label{lem:radial-coherence}
Let $C$ be a reachable cluster and let $z\notin C$.  Then
\[
    \min\{P(z,C),N(z,C)\}\le 2L(z,C).
\]
\end{lemma}

\begin{proof}
The proof is by induction on the merge tree of $C$.  We begin by introducing
normalized notation that will make the induction easier to follow.

For every nonempty set $S$ not containing $z$, define
\begin{align*}
    \widehat P(S)
    &=\frac{P(z,S)}{|S|},\\
    \widehat N(S)
    &=\frac{N(z,S)}{|S|},\\
    \widehat U(S)
    &=\frac{2L(z,S)}{|S|}.
\end{align*}
Thus, the desired inequality for $S$ is equivalent to
\begin{equation}
    \min\{\widehat P(S),\widehat N(S)\}
    \le \widehat U(S).
    \label{eq:normalized-radial-goal}
\end{equation}

It is useful to understand these three quantities probabilistically.  Choose
a vertex $v$ uniformly at random from $S$, and define
\[
    Q=1-2d_{zv}.
\]
Since $0\le d_{zv}\le 1$, we have $Q\in[-1,1]$.  Moreover,
\[
    (Q)_+=(1-2d_{zv})_+=p_{zv},
\]
and
\[
    (-Q)_+=(2d_{zv}-1)_+=n_{zv}.
\]
Also,
\[
    1-|Q|
    =1-|1-2d_{zv}|
    =2\min\{d_{zv},1-d_{zv}\}
    =2\ell_{zv}.
\]
Consequently,
\begin{align}
    \widehat P(S)&=\mathbb E[(Q)_+],\label{eq:normalized-P}\\
    \widehat N(S)&=\mathbb E[(-Q)_+],\label{eq:normalized-N}\\
    \widehat U(S)&=\mathbb E[1-|Q|].\label{eq:normalized-U}
\end{align}
Here the expectation is over the uniformly random choice of $v\in S$.

We say that $S$ admits the \emph{positive orientation} if
\[
    \widehat N(S)\le \widehat U(S),
\]
and that it admits the \emph{negative orientation} if
\[
    \widehat P(S)\le \widehat U(S).
\]
Therefore,~\eqref{eq:normalized-radial-goal} says precisely that $S$ admits
at least one of these two orientations.

\paragraph{Base case.}
Suppose that $C=\{v\}$ is a singleton.  At most one of
$p_{zv}$ and $n_{zv}$ is nonzero.  Hence
\[
    \min\{P(z,C),N(z,C)\}=0\le 2L(z,C).
\]
Thus, the lemma holds for every leaf of the merge tree.

\paragraph{Inductive step.}
Suppose that $C$ was obtained by merging two reachable clusters $A$ and
$B$, and assume inductively that the lemma holds for both children.

The normalized quantities for $C=A\cup B$ are size-weighted averages:
\begin{align*}
    \widehat P(C)
    &=
    \frac{|A|\widehat P(A)+|B|\widehat P(B)}
         {|A|+|B|},\\
    \widehat N(C)
    &=
    \frac{|A|\widehat N(A)+|B|\widehat N(B)}
         {|A|+|B|},\\
    \widehat U(C)
    &=
    \frac{|A|\widehat U(A)+|B|\widehat U(B)}
         {|A|+|B|}.
\end{align*}

If both $A$ and $B$ admit the positive orientation, then
\[
    \widehat N(A)\le\widehat U(A)
    \qquad\text{and}\qquad
    \widehat N(B)\le\widehat U(B).
\]
Taking their size-weighted average gives
\[
    \widehat N(C)\le\widehat U(C).
\]
Hence $C$ also admits the positive orientation.

Similarly, if both children admit the negative orientation, then
\[
    \widehat P(C)\le\widehat U(C).
\]
Therefore, the only potentially difficult situation is when the two children
do not admit any common orientation.

After exchanging the names of $A$ and $B$ if necessary, this difficult case
has the form
\begin{equation}
    \widehat N(A)\le\widehat U(A)<\widehat P(A),
    \qquad
    \widehat P(B)\le\widehat U(B)<\widehat N(B).
    \label{eq:opposite-orientations}
\end{equation}
Indeed, $A$ admits only the positive orientation, while $B$ admits only the
negative orientation.  We shall prove that two clusters satisfying
\eqref{eq:opposite-orientations} cannot satisfy the algorithm's merge rule.

Choose $a$ uniformly from $A$ and $b$ uniformly from $B$, independently.
Thus, every ordered pair $(a,b)\in A\times B$ is selected with probability
$1/(|A||B|)$.  Define the random variables
\[
    X=1-2d_{za},
    \qquad
    Y=2d_{zb}-1.
\]
The sign in the definition of $Y$ is deliberately reversed.  Cluster $A$ is
positively oriented relative to $z$, whereas $B$ is negatively oriented.
The reversal makes both $X$ and $Y$ positively oriented random variables.

For a random variable $Z\in[-1,1]$, define
\[
    P_Z=\mathbb E[(Z)_+],
    \qquad
    N_Z=\mathbb E[(-Z)_+],
    \qquad
    U_Z=\mathbb E[1-|Z|].
\]
Using~\eqref{eq:normalized-P}--\eqref{eq:normalized-U}, we have
\begin{align*}
    P_X&=\widehat P(A),
    &
    N_X&=\widehat N(A),
    &
    U_X&=\widehat U(A),\\
    P_Y&=\widehat N(B),
    &
    N_Y&=\widehat P(B),
    &
    U_Y&=\widehat U(B).
\end{align*}
The exchange of $P$ and $N$ in the second line is caused by the sign reversal
in $Y$.  It follows from~\eqref{eq:opposite-orientations} that
\begin{equation}
    N_X\le U_X<P_X,
    \qquad
    N_Y\le U_Y<P_Y.
    \label{eq:X-Y-orientations}
\end{equation}

We next prove the following scalar claim.

\paragraph{Scalar claim.}
If $X,Y\in[-1,1]$ are independent random variables and both satisfy
\[
    N_Z\le U_Z\le P_Z,
\]
then
\begin{equation}
    \mathbb E|X+Y|\ge\frac12.
    \label{eq:scalar-claim}
\end{equation}

To prove the claim, first observe that, pointwise,
\[
    (Z)_++(-Z)_++1-|Z|=1.
\]
Taking expectations gives
\begin{equation}
    P_Z+N_Z+U_Z=1.
    \label{eq:PNU-sum}
\end{equation}
Therefore, the inequalities $N_Z\le U_Z\le P_Z$ are equivalent to
\[
    P_Z+2N_Z\le1
    \qquad\text{and}\qquad
    2P_Z+N_Z\ge1.
\]
Together with $P_Z,N_Z\ge0$, these inequalities imply that the point
$(P_Z,N_Z)$ belongs to the triangle
\[
    \mathcal T
    =
    \operatorname{conv}
    \left\{
        \left(\frac12,0\right),
        (1,0),
        \left(\frac13,\frac13\right)
    \right\},
\]
where $\operatorname{conv}$ denotes the convex hull.

Let
\[
    \alpha=\Pr[X\ge0],
    \qquad
    \beta=\Pr[Y\ge0].
\]
For every $x,y\in[-1,1]$, we have
\begin{equation}
    |x+y|
    \ge
    (-x)_++(-y)_+
    +x\ind[y\ge0]
    +y\ind[x\ge0].
    \label{eq:pointwise-scalar}
\end{equation}
To see this, consider the signs of $x$ and $y$.

\begin{itemize}
    \item If $x,y\ge0$, then the right-hand side equals $x+y=|x+y|$.
    \item If $x,y<0$, then the right-hand side equals
          $-x-y=|x+y|$.
    \item If $x$ and $y$ have opposite signs, then the right-hand side is
          zero, and hence is at most $|x+y|$.
\end{itemize}

Taking expectations in~\eqref{eq:pointwise-scalar} and using independence,
we obtain
\begin{align}
    \mathbb E|X+Y|
    &\ge
    N_X+N_Y
    +\mathbb E[X]\Pr[Y\ge0]
    +\mathbb E[Y]\Pr[X\ge0]\notag\\
    &=
    N_X+N_Y
    +\beta(P_X-N_X)
    +\alpha(P_Y-N_Y).
    \label{eq:scalar-expectation-one}
\end{align}
Since $N_X\le P_X$ and $N_Y\le P_Y$, both
$P_X-N_X$ and $P_Y-N_Y$ are nonnegative.  Furthermore,
\[
    (X)_+\le\ind[X\ge0]
    \qquad\text{and}\qquad
    (Y)_+\le\ind[Y\ge0],
\]
because $X,Y\le1$.  Therefore,
\[
    \alpha\ge P_X,
    \qquad
    \beta\ge P_Y.
\]
Substituting these lower bounds into
\eqref{eq:scalar-expectation-one} gives
\begin{align}
    \mathbb E|X+Y|
    &\ge
    N_X+N_Y
    +P_Y(P_X-N_X)
    +P_X(P_Y-N_Y)\notag\\
    &=
    N_X+N_Y+2P_XP_Y-P_YN_X-P_XN_Y.
    \label{eq:bilinear-lower-bound}
\end{align}

The right-hand side of~\eqref{eq:bilinear-lower-bound} is bilinear in
$(P_X,N_X)$ and $(P_Y,N_Y)$.  If one pair is fixed, the expression is linear
in the other pair, so its minimum over a triangle occurs at a vertex.
Applying this argument to both pairs shows that it is enough to check the
nine pairs of vertices of $\mathcal T$.  The values are
\[
\begin{array}{c|ccc}
 &(\tfrac12,0)&(1,0)&(\tfrac13,\tfrac13)\\ \hline
(\tfrac12,0)&\tfrac12&1&\tfrac12\\
(1,0)&1&2&\tfrac23\\
(\tfrac13,\tfrac13)&\tfrac12&\tfrac23&\tfrac23
\end{array}
\]
and the smallest value is $1/2$.  This proves the scalar claim.

We now return to the clusters $A$ and $B$.  By the definitions of $X$ and
$Y$,
\begin{align*}
    |X+Y|
    &=
    |1-2d_{za}+2d_{zb}-1|\\
    &=
    2|d_{zb}-d_{za}|.
\end{align*}
The reverse triangle inequality for the semimetric $d$ gives
\[
    |d_{zb}-d_{za}|\le d_{ab}.
\]
Therefore,
\[
    |X+Y|\le2d_{ab}.
\]
Taking expectations over the independent uniform choices of $a$ and $b$,
we obtain
\[
    \mathbb E|X+Y|
    \le
    2\mathbb E[d_{ab}]
    =
    2\frac{D(A,B)}{|A||B|}.
\]
Since the algorithm merged $A$ and $B$, its merge rule gives
\[
    \frac{D(A,B)}{|A||B|}<\frac14.
\]
Consequently,
\[
    \mathbb E|X+Y|<\frac12,
\]
contradicting the scalar claim.

Thus, the difficult opposite-orientation case cannot occur in a valid merge.
The two children must admit a common orientation, and that orientation is
preserved by taking their size-weighted average.  Hence
\[
    \min\{\widehat P(C),\widehat N(C)\}\le\widehat U(C).
\]
Multiplying by $|C|$ gives
\[
    \min\{P(z,C),N(z,C)\}\le2L(z,C),
\]
as required.
\end{proof}
The next lemma converts radial coherence for every row and column of a matrix
into a global inequality.

\begin{lemma}[Polarization matrix lemma]
Let $(q_{ij})$ be a finite matrix with entries in $[-1,1]$.  For every entry,
define
\[
 p_{ij}=(q_{ij})_+,
 \qquad
 n_{ij}=(-q_{ij})_+,
 \qquad
 u_{ij}=1-|q_{ij}|.
\]
For a row or column $\lambda$, let $P_\lambda,N_\lambda,U_\lambda$ denote
the sums of $p_{ij},n_{ij},u_{ij}$ over that row or column.  Suppose that
every row and every column satisfies
\[
  \min\{P_\lambda,N_\lambda\}\le U_\lambda.
\]
If
\[
 P=\sum_{i,j}p_{ij},
 \qquad
 N=\sum_{i,j}n_{ij},
 \qquad
 U=\sum_{i,j}u_{ij},
\]
then
\[
  \min\{P,N\}\le\frac52U.
\]
\end{lemma}

\begin{proof}
Choose an admissible orientation for every row and column.  The orientation
``$+$'' is admissible when $N_\lambda\le U_\lambda$, and ``$-$'' is
admissible when $P_\lambda\le U_\lambda$.  At least one orientation is
admissible; if both are admissible, choose either one.

Let $R_+,R_-$ be the sets of positively and negatively oriented rows, and
let $C_+,C_-$ be the corresponding sets of columns.  Partition the matrix
into four rectangular blocks
\[
 \mathcal A=R_+\times C_+,
 \quad
 \mathcal B=R_+\times C_-,
 \quad
 \mathcal C=R_-\times C_+,
 \quad
 \mathcal D=R_-\times C_-.
\]
Let $a,b,c,d$ be the numbers of entries in these four blocks.  For
$\mathcal X\in\{\mathcal A,\mathcal B,\mathcal C,\mathcal D\}$, define
\[
 P_{\mathcal X}=\sum_{(i,j)\in\mathcal X}p_{ij},
 \quad
 N_{\mathcal X}=\sum_{(i,j)\in\mathcal X}n_{ij},
 \quad
 U_{\mathcal X}=\sum_{(i,j)\in\mathcal X}u_{ij}.
\]
Because $p_{ij}+n_{ij}+u_{ij}=1$, we have
\[
 P_{\mathcal A}+N_{\mathcal A}+U_{\mathcal A}=a,
\]
and analogous identities for the other three blocks.

Put $T=\min\{P,N\}$.  Aggregating the row and column hypotheses gives
\begin{align*}
 g_1&=N_{\mathcal A}+N_{\mathcal B}
      -U_{\mathcal A}-U_{\mathcal B}\le0,\\
 g_2&=P_{\mathcal C}+P_{\mathcal D}
      -U_{\mathcal C}-U_{\mathcal D}\le0,\\
 g_3&=N_{\mathcal A}+N_{\mathcal C}
      -U_{\mathcal A}-U_{\mathcal C}\le0,\\
 g_4&=P_{\mathcal B}+P_{\mathcal D}
      -U_{\mathcal B}-U_{\mathcal D}\le0.
\end{align*}
Furthermore,
\[
 g_5=T-P\le0,
 \qquad
 g_6=T-N\le0.
\]
Direct expansion, using the four block-size identities, gives
\begin{align}
 \frac52U-T
 &=\frac b2-d
   -\left(\frac32g_1+g_3+\frac12g_4+g_6\right)
   +\frac32N_{\mathcal A}+\frac32U_{\mathcal C}
   +\frac32P_{\mathcal D}+3U_{\mathcal D},
 \label{eq:certificate-one}\\
 \frac52U-T
 &=\frac{-a+3b+3c-5d}{6}
   -\left(\frac43g_1+\frac23g_2+\frac43g_3+\frac23g_4
          +\frac16g_5+\frac56g_6\right)
   +2N_{\mathcal A}+2P_{\mathcal D}+2U_{\mathcal D}.
 \label{eq:certificate-two}
\end{align}
Every term following the first term on the right-hand side of either
identity is nonnegative, because $g_1,\ldots,g_6\le0$ and
$P_{\mathcal X},N_{\mathcal X},U_{\mathcal X}\ge0$.

Let $M$ be the total number of matrix entries, and define
\[
 \rho=\frac{|R_+|}{|R_+|+|R_-|},
 \qquad
 \gamma=\frac{|C_+|}{|C_+|+|C_-|}.
\]
Transposing the matrix exchanges $\rho$ and $\gamma$.  Replacing every
$q_{ij}$ by $-q_{ij}$ and reversing every chosen orientation sends
$(\rho,\gamma)$ to $(1-\rho,1-\gamma)$ without changing $T$ or $U$.
Therefore, we may assume
\[
  \rho\ge\gamma,
  \qquad
  \rho+\gamma\ge1.
\]

If $\rho\ge2/3$, then
\[
 \frac b2-d
 =(1-\gamma)\left(\frac{3\rho}{2}-1\right)M\ge0.
\]
Equation~\eqref{eq:certificate-one} now implies $T\le(5/2)U$.

Otherwise, $\rho<2/3$.  Set $s=\rho+\gamma$.  Since
$\rho\ge\gamma$ and $\rho+\gamma\ge1$, we have $1\le s<4/3$.  Also,
\begin{align*}
 \frac{-a+3b+3c-5d}{M}
 &=8s-12\rho\gamma-5\\
 &\ge8s-3s^2-5\\
 &=(s-1)(5-3s)\ge0,
\end{align*}
where we used $\rho\gamma\le s^2/4$.  Equation
~\eqref{eq:certificate-two} again gives $T\le(5/2)U$.
\end{proof}

\begin{corollary}[Pair coherence]
For any two disjoint reachable clusters $A,B$,
\[
  \min\{P(A,B),N(A,B)\}\le5L(A,B).
\]
\end{corollary}

\begin{proof}
Form the $|A|\times|B|$ matrix
\[
  q_{uv}=1-2d_{uv},
  \qquad u\in A,\ v\in B.
\]
Fixing $u\in A$ and applying radial coherence to the reachable cluster $B$
shows that the corresponding row satisfies the hypothesis of the
polarization matrix lemma.  Fixing $v\in B$ and applying radial coherence to
$A$ proves the same for every column.  The total matrix ambiguity is
\[
 \sum_{u\in A}\sum_{v\in B}(1-|q_{uv}|)
 =2L(A,B).
\]
The polarization matrix lemma therefore gives
\[
 \min\{P(A,B),N(A,B)\}
 \le\frac52\cdot2L(A,B)=5L(A,B).
\]
\end{proof}

\begin{proof}[Proof of the main theorem]
Fix a time $t$.  Consider an internal node of the merge tree of an output
cluster, and let $A,B$ be its two children.  When they were merged,
\[
  D(A,B)<\frac{|A||B|}{4}.
\]
By~\eqref{eq:set-identities},
\[
 P(A,B)-N(A,B)
 =|A||B|-2D(A,B)
 >\frac{|A||B|}{2}>0.
\]
Thus $N(A,B)$ is the smaller of $P(A,B)$ and $N(A,B)$.  Pair coherence gives
$N(A,B)\le5L(A,B)$, and the cost paid by joining all pairs in $A\times B$ is
\[
 D(A,B)=L(A,B)+N(A,B)\le6L(A,B).
\]
Every pair of vertices lying in the same output cluster appears in exactly
one such merge-tree cut: the cut corresponding to the lowest common ancestor
of its two leaves.

Now let $A,B$ be two distinct output clusters at the end of time $t$.  Since
the merge loop has terminated,
\[
  D(A,B)\ge\frac{|A||B|}{4}.
\]
The algorithm separates all pairs in $A\times B$, so its cost on these pairs
is
\[
 |A||B|-D(A,B)=L(A,B)+P(A,B).
\]
If $P(A,B)\le N(A,B)$, pair coherence gives $P(A,B)\le5L(A,B)$, and hence
the separated cost is at most $6L(A,B)$.

Otherwise, $N(A,B)<P(A,B)$.  Pair coherence gives
$N(A,B)\le5L(A,B)$, and therefore
\[
  D(A,B)=L(A,B)+N(A,B)\le6L(A,B).
\]
Terminality implies $|A||B|\le4D(A,B)$, and hence
\[
 |A||B|-D(A,B)
 \le3D(A,B)
 \le18L(A,B).
\]

The internal merge-tree cuts together with the pairs of distinct terminal
clusters partition all unordered pairs of $V_t$.  Summing the preceding
bounds and using~\eqref{eq:ambiguity-lower-bound}, we obtain
\[
 \cost_t(\mathcal C_t)
 \le18L_t
 \le18\OPT_t.
\]

It remains to bound recourse.  A vertex changes its cluster identifier only
when its current cluster is the smaller child, or the non-inheriting child in
a tie, of a merge.  If that child's size is $s$, the merged cluster has size
at least $2s$.  Thus, whenever a vertex incurs one unit of recourse, the size
of its current cluster at least doubles.  Since cluster sizes are at most
$n$, every vertex incurs at most $\lfloor\log_2 n\rfloor$ units of recourse.
This bound contains no dependence on $m$.
\end{proof}

\begin{remark}[Metric weighted correlation clustering]
The analysis uses only that $d_{uv}\in[0,1]$ and that $d$ satisfies the
triangle inequality.  Therefore, the same algorithm and proof apply to the
online metric weighted correlation clustering model in which joining
$u,v$ costs $d_{uv}$ and separating them costs $1-d_{uv}$.
\end{remark}

\end{document}